\section{Conclusion}
\label{sec:conclusion}

This work reviewed recent applications of large language models in safety-critical human–machine interaction, encompassing text-based, audio-based, and visual language interfaces. 
Across the examined domains, LLMs demonstrate the greatest value when deployed as assistive systems that support interpretation, planning, and decision-making rather than as autonomous controllers. 
Studies in aerospace, emergency response, and industrial settings consistently highlight the effectiveness of hybrid architectures that separate high-level reasoning from deterministic low-level execution.

Overall, the findings indicate that successful integration of LLMs, including speech-enabled and vision-language models, depends on constrained agentic designs, human-in-the-loop control, and clear system boundaries. 
While challenges related to bias, calibration, multimodal integration, and reliability remain, the evidence suggests that carefully designed multimodal LLM-based interfaces can reduce cognitive load and improve usability without compromising safety. 
Future work should focus on retrieval-grounded reasoning, robust multimodal fusion, and human-centered evaluation to support deployment in real-world safety-critical environments.
